\documentclass[10pt,a4paper]{extarticle}
\usepackage{parskip}

\usepackage[utf8]{inputenc}
\usepackage[T1]{fontenc}
\usepackage[brazil]{babel}
\usepackage{amsmath, amssymb, amsthm}
\usepackage{geometry}
\usepackage{xcolor}
\usepackage{hyperref}
\usepackage{booktabs}
\usepackage{array}
\usepackage{tabularx}
\usepackage{enumitem}
\usepackage[protrusion=true,expansion=false]{microtype}
\usepackage{csquotes}
\usepackage{tcolorbox}

\definecolor{important}{RGB}{255, 100, 100}
\definecolor{notice}{RGB}{0, 102, 204}
\definecolor{accent}{RGB}{230, 120, 30}

\newcommand{\impt}[1]{\textcolor{important}{#1}}
\newcommand{\ntc}[1]{\textcolor{notice}{#1}}

\setlist{nosep}
\geometry{margin=2cm}

\newtcolorbox{respbox}{
  colback=notice!5,
  colframe=notice!50,
  boxrule=0.4pt,
  arc=2pt,
  left=4pt, right=4pt, top=3pt, bottom=3pt,
  fontupper=\small
}

\title{Lista de Exercícios -- Tipos de Pesquisa\\
\large Tipos de Pesquisa em Ciência da Computação\footnote{Baseado em Wazlawick, R. S. \textit{Metodologia de Pesquisa para Ciência da Computação}, 3.ed. Elsevier, 2020.}}
\author{BCC502 -- Metodologia Científica para Ciência da Computação}
\date{\today}

\begin{document}
\maketitle


\begin{enumerate}[label=\textbf{\arabic*.}]

\item Apresente  as definições dos tipos de pesquisa abaixo de acordo com \texttt{Wazlawick, R. S. \textit{Metodologia de Pesquisa para Ciência da Computação}, 3.ed. Elsevier, 2020.} 
    
\begin{enumerate}[label=(\alph*)]
    \item Apresentação de um Produto
    \item Apresentação de Algo Diferente
    \item Apresentação de Algo Presumivelmente Melhor 
    \item Apresentação de Algo Reconhecidamente Melhor
    \item Apresentação de uma Prova
\end{enumerate}

  \item  Escolha um problema da sua área de interesse (ex.: aprendizado de máquina, segurança, engenharia de \textit{software}, sistemas, IHC, etc.). Para esse mesmo problema, proponha três projetos distintos, cada um pertencente a um tipo de contribuição diferente do exercício anterior. Em cada proposta, explicite:
  \begin{enumerate}[label=(\alph*)]
    \item A questão de pesquisa.
    \item O tipo de contribuição.
    \item O procedimento técnico predominante.
    \item O critério que permitirá afirmar que o projeto produziu conhecimento novo.
  \end{enumerate}
\end{enumerate}

\end{document}